% chapters/03/section-06.tex
\section{Menjalankan, Memeriksa, dan Menelusuri Kesalahan}
\label{sec:menelusuri-kesalahan}

Kemampuan menelusuri kesalahan sama pentingnya dengan kemampuan menulis kode.
Bagian ini melatih dua hal yang akan Anda pakai setiap hari: menjalankan
program dan membaca pesan kesalahan.

\subsection{Alur Kerja Harian}

Menulis program bukan kegiatan sekali jadi. Pekerjaan berlangsung sebagai
perulangan pendek: tulis, jalankan, amati, perbaiki bila perlu, lalu simpan
hasilnya. Perulangan tersebut diperlihatkan pada Gambar~\ref{fig:alur-kerja}.

\begin{figure}[htbp]
\centering
% chapters/03/figures/alur-kerja.tikz
% Alur kerja harian: tulis, jalankan, amati, simpan; kembali menulis bila
% hasilnya belum sesuai harapan.
\begin{tikzpicture}[
  langkah/.style={
    rectangle, rounded corners=2pt, draw=ubPrimary, thick, fill=ubAlt,
    align=center, text width=26mm, inner sep=3.5pt, minimum height=11mm,
    font=\footnotesize
  },
  akhir/.style={langkah, draw=ubGreen, fill=ubGreenBg},
  panah/.style={-{Stealth[length=2.4mm]}, thick, draw=ubAccent},
  ulang/.style={-{Stealth[length=2.4mm]}, thick, draw=ubWarn, dashed},
  label/.style={font=\scriptsize, ubWarn}
]

\node[langkah] (tulis) at (0,0) {Tulis atau\\perbaiki kode};
\node[langkah, right=8mm of tulis] (jalan) {Jalankan\\program};
\node[langkah, right=8mm of jalan] (amati) {Amati keluaran\\dan pesan galat};
\node[akhir, right=8mm of amati] (simpan) {Simpan dan\\catat hasilnya};

\draw[panah] (tulis) -- (jalan);
\draw[panah] (jalan) -- (amati);
\draw[panah] (amati) -- (simpan);

\draw[ulang] (amati.south) -- ++(0,-10mm) -| (tulis.south);
\node[label, anchor=north] at (4.5,-1.5)
  {bila hasilnya belum sesuai harapan};
\end{tikzpicture}

\caption{Alur kerja harian dari menulis kode sampai menyimpan hasil}
\label{fig:alur-kerja}
\end{figure}

Perulangan itu sebaiknya dipendekkan sebanyak mungkin. Perubahan kecil yang
langsung dijalankan lebih mudah dilacak daripada perubahan besar yang
melibatkan puluhan baris sekaligus, karena apabila terjadi kesalahan Anda
masih mengingat bagian mana yang baru saja disentuh.

\subsection{Cara Menjalankan Program}

Program dijalankan dari terminal, tepatnya dari folder tempat berkas program
berada.

\begin{kodekeluaran}
cd proyek-pygame\code
python ch03_jendela.py
\end{kodekeluaran}

Untuk mengubah ukuran jendela tanpa menyunting kode, tambahkan pilihan pada
perintah.

\begin{kodekeluaran}
python ch03_jendela.py --width 800 --height 600
\end{kodekeluaran}

Apabila program dijalankan melalui tombol pada editor, pastikan editor
menjalankan berkas dengan penafsir Python yang sama dengan yang dipakai pada
pemeriksaan lingkungan. Perbedaan penafsir merupakan penyebab paling sering
munculnya pesan bahwa Pygame tidak ditemukan, walaupun Pygame sudah terpasang.

\subsection{Membaca Pesan Kesalahan}

Pesan kesalahan Python berupa \istilah{traceback}, yaitu jejak urutan
pemanggilan yang berakhir pada baris penyebab kesalahan. Kebiasaan yang paling
menghemat waktu adalah \textbf{membaca dari baris terakhir ke atas}, bukan dari
atas ke bawah.

Perhatikan dua contoh nyata berikut.

\begin{kodekeluaran}
File "proyek-pygame\code\galat1.py", line 3
    screen = pygame.display.set_mode((640, 480)
                                    ^
SyntaxError: '(' was never closed
\end{kodekeluaran}

Pesan tersebut memuat tiga keterangan yang cukup untuk menemukan penyebabnya:
jenis kesalahan (\code{SyntaxError}), nomor baris (3), dan penunjuk berupa tanda
\code{\^{}} pada posisi yang bermasalah. Di sini kurung pembuka
\code{set_mode(} tidak pernah ditutup.

\begin{kodekeluaran}
Traceback (most recent call last):
  File "proyek-pygame\code\galat2.py", line 5, in <module>
    screen.fill(WARNA_LATAR)
                ^^^^^^^^^^^
NameError: name 'WARNA_LATAR' is not defined
\end{kodekeluaran}

Baris terakhir menyatakan bahwa nama \code{WARNA_LATAR} belum dikenal.
Penyebabnya ada dua kemungkinan: nama itu belum pernah dibuat, atau dibuat
dengan ejaan yang berbeda. Pada pemeriksaan lebih lanjut, kemungkinan kedua
yang berlaku---nama yang benar adalah \code{WARNA_LATAR} dengan huruf besar
seluruhnya, sedangkan pada baris sebelumnya tertulis dengan huruf kecil.

Pada contoh kedua, program sudah sempat berjalan sampai baris
kelima. Kesalahan nama seperti itu tidak dapat dideteksi sebelum barisnya
dieksekusi. Karena itu, cara paling andal menemukan kesalahan adalah
menjalankan programnya, bukan membacanya berulang kali.

\subsection{Peringatan Bukan Kesalahan}

Terminal sering menampilkan pesan yang tampak mengkhawatirkan tetapi bukan
kesalahan. Perbedaan keduanya penting, karena mengejar peringatan yang tidak
perlu hanya membuang waktu.

\begin{description}
  \item[Kesalahan (\istilah{error})] Menghentikan program, atau membuat
        program tidak dapat dijalankan sama sekali. Selama masih ada pesan
        kesalahan, program tidak akan berjalan.
  \item[Peringatan (\istilah{warning})] Memberitahukan sesuatu yang sebaiknya
        diketahui, tetapi program tetap berjalan. Peringatan boleh diabaikan
        apabila Anda memahami maksudnya.
\end{description}

Ketika program jendela pertama dijalankan, muncul tiga pesan yang semuanya
termasuk peringatan, bukan kesalahan.

\begin{kodekeluaran}
pygame 2.6.1 (SDL 2.28.4, Python 3.11.9)
Hello from the pygame community. https://www.pygame.org/contribute.html
.../pygame/pkgdata.py:25: UserWarning: pkg_resources is deprecated ...
\end{kodekeluaran}

Dua baris pertama adalah sambutan Pygame yang menampilkan versi beserta ajakan
mengunjungi laman komunitas. Baris ketiga adalah peringatan bahwa Pygame masih
memakai bagian pustaka yang sudah usang. Ketiganya berasal dari Pygame, bukan
dari program Anda, dan tidak memengaruhi jalannya program.

Yang menandakan berjalannya program dengan baik adalah keluaran dari program
Anda sendiri, yaitu baris seperti \code{Jendela 640x480 dibuka.} dan
\code{Program selesai.} Sesudah jendela ditutup. Apabila kedua baris itu
muncul, program Anda bekerja sebagaimana mestinya.

\subsection{Kesalahan yang Sering Terjadi}

Tabel~\ref{tab:galat-umum} memuat kesalahan yang paling sering muncul pada
tahap ini beserta cara menanganinya. Simpan tabel ini; hampir semua masalah
yang Anda hadapi pada minggu-minggu pertama ada di dalamnya.

\begin{table}[htbp]
\centering
\caption{Kesalahan yang sering terjadi pada tahap penyiapan dan penanganannya}
\label{tab:galat-umum}
{\setlength{\tabcolsep}{4pt}\footnotesize%
\begin{tabularx}{\textwidth}{@{}L{4.4cm} L{4.0cm} Y@{}}
\toprule
\headrow \textbf{Pesan atau gejala} & \textbf{Penyebab} & \textbf{Perbaikan} \\
\midrule
\code{SyntaxError: '(' was never closed} &
Ada kurung, tanda kutip, atau tanda kurung siku yang belum ditutup. &
Periksa baris yang ditunjuk tanda panah; hitung pasangan kurung buka dan
tutupnya. \\
\addlinespace
\code{IndentationError: unexpected indent} &
Indentasi tidak konsisten, misalnya bercampur tab dan spasi. &
Pakai empat spasi untuk setiap tingkat. Atur editor agar mengubah tab menjadi
spasi secara otomatis. \\
\addlinespace
\code{NameError: name 'X' is not defined} &
Nama dipakai sebelum dibuat, atau ejaannya berbeda dengan yang dibuat. &
Periksa ejaan dan perhatikan besar-kecil huruf; pastikan baris pembuatan
variabel berada lebih dahulu. \\
\addlinespace
\code{ModuleNotFoundError: No module named 'pygame'} &
Pygame belum terpasang pada penafsir yang sedang dipakai. &
Jalankan skrip pemeriksaan lingkungan, lalu pasang dengan
\code{python -m pip install pygame}. \\
\addlinespace
\code{pygame.error: No available video device} &
Program dijalankan pada lingkungan tanpa layar. &
Kerjakan pada komputer laboratorium; untuk pengujian tanpa layar, setel
\code{SDL_VIDEODRIVER} bernilai \code{dummy}. \\
\addlinespace
\code{FileNotFoundError: [Errno 2] No such file or directory} &
Folder kerja terminal bukan folder tempat berkas berada. &
Jalankan \code{cd} ke folder yang benar, lalu periksa isinya dengan \code{dir}
atau \code{ls}. \\
\bottomrule
\end{tabularx}}
\end{table}
